\section{Zadanie 4}
Załóżmy, że mamy do dyspozycji generator liczb z rozkładu jednostajnego z przedziału \([-1, 1]\).
Naszym celem będzie wygenerowanie liczb z rozkładu normalnego \(N(0,1)\).
Niech \(U,V\) będą wygenerowanymi niezależnie wartościami na \([-1, 1]\).
Procedura wygląda następująco. Jeśli \(U^2+V^2 > 1\), to odrzucamy wynik i próbujemy jeszcze raz.
W przeciwnym przypadku niech \(S = U^2 + V^2\) i niech
\[X = U\sqrt{\frac{-2 \text{ln}S} {S}}, Y = V\sqrt{\frac{-2 \text{ln}S} {S}}.\]
Udowodnij, że \(X,Y\) są niezależnymi zmiennymi losowymi z rozkładu \(N(0,1)\).
Pamiętaj o wstępnym warunkowaniu! Wskazówka: przeczytaj sekcję \(9.5\) w podręczniku.

\subsection*{Rozwiązanie}
Warunek \(U^2+V^2 \le 1\), oznacza że losujemy punktu z koła z równym prawdopodobieństwem.
Wtedy niech \(\frac{U}{\sqrt{S}} = \)cos\(\alpha\), \(\frac{V}{\sqrt{S}} = \)sin\(\alpha\), z tego mamy:
\[X = \text{cos}\alpha\sqrt{-2\text{ln}S}, Y = \text{sin}\alpha\sqrt{-2\text{ln}S}\]
Z tego że losujemy punkty z koła, mamy że \(\alpha\) jest niezależna od \(S\) oraz ma rozkład jednostajny w \([0, 2\pi]\).
\[P(S < s) = P(U^2 + V^2 < s) = \frac{\pi s}{\pi} = s\]
Z tego wynika że \(S\) ma rozkład jednostajny w \([0, 1]\).

Niech \(X^{\prime}, Y^{\prime}\) - będą niezależnymi zmiennymi z rozkładem \(N(0,1)\), wtedy:
\[F(x^\prime, y^\prime) = \int_{-\infty}^{y^\prime}\int_{-infty}^{x^\prime}\frac{1}{2\pi}e^{-(x^2 + y^2)/2}dxdy\]

Zauważymy że wybieramy punkt na płaszczyźnie, to przejdziemy na współrzędne biegunowe z \(R^2 = {X^{\prime}}^2 + {Y^{\prime}}^2\), a \(\theta = \tan^{-1}\frac{X^\prime}{Y^\prime}\).
Wtedy:
\[ X^\prime = R\cos\theta, Y^\prime = R\sin\theta\].
Za pomocom Jacobiana otrzymujemy:
\[F(r^\prime, \theta^\prime) = \int_{0}^{\theta^\prime}\int_{0}^{r^\prime}\frac{1}{2\pi}e^{-r^2/2}rdrd\theta = \frac{\theta^\prime}{2\pi}(1 - e^{-{r^\prime}^2/2})\]
Zauważymy że \(F(r^\prime, \theta^\prime) = G(r^\prime)H(\theta^\prime)\), gdzie \(G(r^\prime) = 1 - e^{-{r^\prime}^2/2}\), \(H(\theta^\prime) = \frac{\theta^\prime}{2\pi}\), to oznacza że \(R\) oraz \(\theta\) są niezależne.
\[ P(R \ge r^\prime) = e^{-{r^\prime}^2/2}, P(S \ge s) = 1 - s\]
Podstawimy \(\theta = \alpha, R = \sqrt{-2\text{ln}S}\), wtedy:
\[X^\prime = \text{cos}\alpha\sqrt{-2\text{ln}S}, Y^\prime = \text{sin}\alpha\sqrt{-2\text{ln}S} \implies X^\prime = X, Y^\prime = Y\]

